\paragraph{Comparing quantum readouts at fixed energy and resolution.}
The common neutral observable algebra permits a precise comparison with a
declared effective field model. A fixed ultraviolet regulator is allowed;
the statement below does not require an interacting limit at arbitrarily
high energy. It separates a generator estimate, a state map, and a detector
estimate. None of those comparison estimates follows from a classical
invariant sector alone.

Let $H_{\rm e}\geq0$ and $H_h$ be self-adjoint operators on separable
complex Hilbert spaces $\mathcal H_{\rm e}$ and $\mathcal H_h$. Write
$U_j(t)=\exp(-itH_j/\hbar)$, with $\hbar>0$, and let
$J:\mathcal H_{\rm e}\to\mathcal H_h$ be a linear isometry. Fix $E>0$
and the \emph{fixed spectral projection}
$P=\mathbf1_{[0,E]}(H_{\rm e})$. Thus $P$ commutes with $U_{\rm e}(t)$
and $H_{\rm e}P$ is bounded. A neutral comparison uses neutral Hilbert
spaces and effects throughout.

\begin{theorem}[Energy-window effect comparison]
\label{thm:effective-quantum-comparison}
Suppose
\begin{equation}
 J\operatorname{Ran}P\subset\mathcal D(H_h),\qquad
 \|(H_hJ-JH_{\rm e})P\|\leq\epsilon_H.
 \label{eq:effective-quantum-full-residual}
\end{equation}
The residual is a bounded map into the \emph{full} target Hilbert space;
its component orthogonal to $J\operatorname{Ran}P$ is included. Let
$0\leq F_h\leq I$ and $0\leq F_{\rm e}\leq I$ be bounded effects with
\begin{equation}
 \|P(J^*F_hJ-F_{\rm e})P\|\leq\epsilon_D.
 \label{eq:effective-quantum-detector-error}
\end{equation}
Let $\rho$ and $\rho_h$ be density operators, and put
$D(\alpha,\beta)=\tfrac12\|\alpha-\beta\|_1$. Assume
\[
 D(\rho_h,J\rho J^*)\leq\epsilon_0,\qquad
 \eta_\rho:=\operatorname{Tr}((I-P)\rho)\leq\eta<1,
 \qquad \epsilon_0,\epsilon_H,\epsilon_D,\eta\geq0.
\]
For every real $t$, the probabilities
\[
 p_h(t)=\operatorname{Tr}\bigl(F_hU_h(t)\rho_hU_h(t)^*\bigr),\qquad
 p_{\rm e}(t)=\operatorname{Tr}\bigl(F_{\rm e}U_{\rm e}(t)
                                      \rho U_{\rm e}(t)^*\bigr)
\]
obey
\begin{equation}
 |p_h(t)-p_{\rm e}(t)|
 \leq\min\left\{1,\ \epsilon_0+2\sqrt\eta+
                  \frac{|t|\epsilon_H}{\hbar}+\epsilon_D\right\}.
 \label{eq:effective-quantum-probability-error}
\end{equation}
In particular, an exactly energy-supported preparation has no tail term.
The same bound is uniform over any declared detector class satisfying
\eqref{eq:effective-quantum-detector-error} with one constant.
\end{theorem}

\begin{proof}
Put $\mathcal K=\operatorname{Ran}P$ and
$A=H_{\rm e}|_{\mathcal K}$. The spectral theorem gives
$\|A\|\leq E$. For $J_E=J|_{\mathcal K}$, the residual
$R=H_hJ_E-J_EA$ is bounded by hypothesis. Consequently
$H_hJ_E=J_EA+R$ is bounded from $\mathcal K$ to $\mathcal H_h$;
in particular $J_E$ is continuous into the operator domain equipped with
its graph norm. For $x\in\mathcal K$ and fixed $t\geq0$, the path
\[
 s\longmapsto U_h(t-s)J_Ee^{-isA/\hbar}x
\]
is strongly differentiable on $[0,t]$, with derivative
\[
 \frac{i}{\hbar}U_h(t-s)R e^{-isA/\hbar}x.
\]
Indeed, the bounded $A$ orbit is differentiable; its image lies in
$\mathcal D(H_h)$, and its Hamiltonian image is continuous by the
displayed bounded expression for $H_hJ_E$. The usual difference quotient
for the product therefore uses the strong derivative of $U_h$ on its
operator domain. This verifies the domain step for an unbounded $H_h$;
it does not assume operator-norm differentiability of $U_h$.
Integration and unitarity give
\begin{equation}
 \|(U_h(t)J-JU_{\rm e}(t))P\|
       \leq |t|\epsilon_H/\hbar.
 \label{eq:effective-quantum-duhamel}
\end{equation}
The same argument on the reversed interval proves the result for $t<0$.

Let $p=\operatorname{Tr}(P\rho)=1-\eta_\rho>0$ and
$\rho_E=P\rho P/p$. A purification $\Psi$ of $\rho$ has normalized
projected purification
$\Psi_E=(P\otimes I)\Psi/\sqrt p$ with overlap $\sqrt p$.
For unit vectors $x,y$, diagonalization on their two-dimensional span gives
\[
 D(|x\rangle\langle x|,|y\rangle\langle y|)
       =\sqrt{1-|\langle x,y\rangle|^2}\leq\|x-y\|.
\]
Contractivity under partial trace proves the normalized conditioning bound
\begin{equation}
 D(\rho,\rho_E)\leq\sqrt{\eta_\rho}\leq\sqrt\eta.
 \label{eq:effective-quantum-conditioning}
\end{equation}
This argument retains coherences across the energy cut; no commutation of
$\rho$ with $P$ is assumed.

Apply \eqref{eq:effective-quantum-duhamel} to a purification of $\rho_E$.
Tensoring a bounded Hilbert-space operator with the identity preserves its
norm. The two evolved purifications therefore differ by at most
$|t|\epsilon_H/\hbar$, and their reduced states have at most this trace
distance. Initial preparation contributes $\epsilon_0$ by unitary
invariance. Replacing $\rho$ by $\rho_E$ on each side contributes at most
$2\sqrt\eta$ by \eqref{eq:effective-quantum-conditioning}, isometric
embedding, and unitary invariance.

Finally $U_{\rm e}(t)\rho_EU_{\rm e}(t)^*$ is supported in $P$.
Its expectations of $J^*F_hJ$ and $F_{\rm e}$ differ by at most
$\epsilon_D$. For any two density operators and any effect $0\leq F\leq I$,
$|\operatorname{Tr}F(\alpha-\beta)|\leq D(\alpha,\beta)$: the positive
and negative parts of the trace-zero self-adjoint difference each have
trace $\|\alpha-\beta\|_1/2$. These facts and the triangle inequality give
the second bound in \eqref{eq:effective-quantum-probability-error}; both
probabilities lie in $[0,1]$, giving the first.
\end{proof}

If a preparation has the form-energy bound
$\operatorname{Tr}(H_{\rm e}^{1/2}\rho H_{\rm e}^{1/2})\leq\overline E$,
then $\eta_\rho\leq\overline E/E$. This follows by integrating
$\mathbf1_{(E,\infty)}(\lambda)\leq\lambda/E$ against its spectral
probability measure. A finite $s$th spectral moment similarly gives
$\eta_\rho\leq M_s/E^s$ for $s>0$. Such tail bounds require their own
certificates; a Gaussian label or a spatial bandwidth alone supplies none.

\paragraph{Why compression is insufficient.}
Take $\mathcal H_{\rm e}=\mathbb C$, $H_{\rm e}=0$,
$\mathcal H_h=\mathbb C^2$, $Jz=(z,0)$, and
$H_h=\left(\begin{smallmatrix}0&1\\1&0\end{smallmatrix}\right)$.
The compressed residual $J^*H_hJ-H_{\rm e}$ vanishes. The full residual
$H_hJ$ has norm one. With $F_{\rm e}=I$, $F_h=JJ^*$ and the exactly
mapped initial pure state, the probabilities are $1$ and
$\cos^2(t/\hbar)$. They differ by one at $t=\pi\hbar/2$.
Thus a compressed Galerkin matrix can pass an exact algebraic comparison
while its omitted coupling changes the observed probability completely.
Lean verifies the interaction-picture norm estimate, the constant-one
bounded-effect inequality and this explicit two-level control. The
spectral-domain and mixed-state trace-class steps above are analytic proofs.

For the charged Whitney Hamiltonian, the self-adjoint neutral operator and
bounded neutral detector functions are available from
Theorems~\ref{thm:whitney-interacting-quantum} and
\ref{thm:whitney-common-observables}. Applying this comparison additionally
requires an effective Hamiltonian, an isometry respecting the intended
preparations, and certified residual, detector and tail bounds on that same
pair of models. The full curved kinetic operator enters the residual.
Neither its coefficients nor couplings to omitted modes may be discarded.
The theorem supplies no numerical comparison accuracy, quantum refinement
limit, experimental identification or source-selected action.
